% Einheit 5 von 5 – Verteilung im Diagramm (bank/binomialverteilung/e5.jsonl, zusammenbau v0.1)
%% TODO Zeitmarke Einheit 5: Marken-Zeile nicht lesbar
%% TODO Zweigzeile Teil 1: was hier gelernt wird (Satz in Schülersprache aus dem Einheitstitel) – Einheit 5
\einheitenkopf[e5]{Einheit 5 von 5 · Verteilung im Diagramm}
\zweigzeile{Abitur GK}
\setcounter{aufgabe}{21}

%% TODO Titel: Ich-kann-Satz fehlt in der Bank (Platzhalter: Kettenname) – Nr. 22
%% TODO Anweisung: ein Satz über der ersten Teilaufgabe fehlt in der Bank – Nr. 22
\begin{aufgabe}{Verteilung im Diagramm}
%% TODO \rechenplatz{n}: Schrittzahl des Grundfalls fehlt in der Bank (nur bei fester Schreibform, 2.3 b)
\begin{teile}
\teil Ein Elfmeterschütze trifft mit 80\,\%. Zu 10 Schüssen gehört der Term $(10 \text{ über } 3) \cdot 0{,}2^3 \cdot 0{,}8^7$. Was wird mit der 3 gezählt? \kreuz{Tore} \\ \kreuz{Fehlschüsse}
\teil Das Säulendiagramm zeigt die Verteilung der Anzahl X der Treffer (Werte gerundet). Lies $P(X = 2)$ ab.

\saeulenab[ymax=35,ystep=5,yfein=1,ylabel=$P(X = k)$ in \%]{0/8,1/24,2/33,3/24,4/10,5/2,6/0}
\leerfeld
\teil 15\,\% der Äpfel einer Plantage haben Druckstellen. 90 Äpfel werden zufällig ausgewählt; X ist die Anzahl mit Druckstellen. Welche Anzahl tritt mit der größten Wahrscheinlichkeit auf? \hfill (Abitur 2017 GK) \\ \leerfeld
\teil Ein Würfelspiel gelingt mit der Wahrscheinlichkeit $0{,}3$ und wird 40-mal gespielt. Ben sagt: „Am wahrscheinlichsten gelingt es 20-mal, das ist die Hälfte.“ Beurteile ohne Wahrscheinlichkeiten zu berechnen.
\teil X ist binomialverteilt mit 10 Versuchen und der Trefferwahrscheinlichkeit $0{,}3$. Genau eine der drei Abbildungen zeigt die Verteilung von X (Werte in Prozent). Entscheide begründet, welche \hfill (Abitur 2023 GK).

\textbf{Abb. 1} \saeulenab[ymax=50,ystep=10,yfein=2,ylabel=\%]{0/0,1/1,2/4,3/12,4/21,5/25,6/21,7/12,8/4,9/1,10/0} \quad \textbf{Abb. 2} \saeulenab[ymax=50,ystep=10,yfein=2,ylabel=\%]{0/3,1/12,2/23,3/27,4/20,5/10,6/4,7/1,8/0,9/0,10/0} \quad \textbf{Abb. 3} \saeulenab[ymax=50,ystep=10,yfein=2,ylabel=\%]{0/5,1/22,2/41,3/49,4/36,5/18,6/7,7/2,8/0,9/0,10/0}
\teil Das Säulendiagramm zeigt die Verteilung der Anzahl X der Treffer bei 8 Versuchen (Werte gerundet). Y = 8 − X zählt die Nieten. Lies $P(Y = 6)$ ab.

\saeulenab[ymax=30,ystep=5,yfein=1,ylabel=$P(X = k)$ in \%]{0/3,1/14,2/26,3/28,4/19,5/8,6/2,7/0,8/0}
\leerfeld
\teil Das Diagramm zeigt die kumulierten Wahrscheinlichkeiten $P(X \le k)$ einer Binomialverteilung mit 6 Versuchen (Werte gerundet). Lies $P(X \le 3)$ ab.

\saeulenab[ymax=100,ystep=20,yfein=2,ylabel=$P(X \le k)$ in \%]{0/3,1/16,2/44,3/74,4/93,5/99,6/100}
\leerfeld
\teil X ist binomialverteilt mit 12 Versuchen und der Trefferwahrscheinlichkeit $0{,}5$; es gilt $P(X \le 4) \approx 0{,}1938$. Bestimme ohne Rechner $P(5 \le X \le 7)$. \leerfeld
\teil X ist binomialverteilt mit 15 Versuchen und der Trefferwahrscheinlichkeit $0{,}5$; die Verteilung ist symmetrisch. Es gilt $P(X \ge 6) \approx 0{,}85$ und $P(X = 9) \approx 0{,}15$. Bestimme damit einen Näherungswert für $P(X = 7)$ \hfill (Abitur 2025 GK). \leerfeld
\end{teile}
\end{aufgabe}

%% TODO Titel: Ich-kann-Satz fehlt in der Bank (Platzhalter: Kettenname) – Nr. 23
%% TODO Anweisung: ein Satz über der ersten Teilaufgabe fehlt in der Bank – Nr. 23
\begin{aufgabe}{Obere Grenze einer im Diagramm markierten kumulierten Wahrscheinlichkeit über den Erwartungswert ermitteln}
\begin{teile}
\teil X ist binomialverteilt mit 24 Versuchen und der Trefferwahrscheinlichkeit $0{,}5$. In einem Diagramm sind die Säulen von 0 bis zu einer Grenze k markiert; zusammen ergeben sie etwa 58\,\%. Bestimme k über den Erwartungswert. k = \leerfeld
\end{teile}
\end{aufgabe}

%% TODO Titel: Ich-kann-Satz fehlt in der Bank (Platzhalter: Kettenname) – Nr. 24
%% TODO Anweisung: ein Satz über der ersten Teilaufgabe fehlt in der Bank – Nr. 24
\begin{aufgabe}{Verteilung im Diagramm – Fehler finden, Begründen, Anwendung, Darstellungswechsel}
\begin{teile}
\teil X ist binomialverteilt mit 40 Versuchen und der Trefferwahrscheinlichkeit $0{,}2$. Lars sagt: „Die höchste Säule liegt in der Mitte, bei 20.“ Finde den Fehler.
\teil Begründe, warum sich die Säulenhöhen im Diagramm einer Binomialverteilung zu eins addieren müssen.
\teil Eine Bäckerei backt täglich 200 Brötchen; jedes misslingt unabhängig mit 3\,\%. Welche Anzahl misslungener Brötchen ist an einem Tag am wahrscheinlichsten?
\teil Eine faire Münze wird viermal geworfen; X ist die Anzahl der Würfe mit Zahl. Zeichne das Säulendiagramm der Verteilung von X (Werte in Prozent).

\saeulen{0/,1/,2/,3/,4/}{40}{10}{$P(X = k)$ in \%}
\end{teile}
\end{aufgabe}
